\chapter[Sermon 66. The Judgment of the Lord Is Described under the Parables of Harvest and Vintage.]{The Judgment of the Lord Is Described under the Parables of Harvest and Vintage.}
\label{sermon:066}
\markright{Sermon 66. The Judgment of the Lord Is Described under the ...}

And I looked, and behold, a white cloud, and upon the cloud one sitting, like unto the Son of man, having on his head a golden crown, and in his hand a sharp sickle. And another angel came out of the Temple, crying with a loud voice to him that sat on the cloud. Thrust in the sickle and reap, for the corn of the Earth is ripe. And he that sat on the cloud thrust in his sickle on the Earth, and the Earth was reaped. And another angel came out of the Temple, which is in Heaven, having also a sharp sickle. And another angel came out from the Altar, which had power over fire, and cried with a loud voice unto him that had the sharp sickle, and said: Thrust in thy sharp sickle, and gather the clusters of the Earth, for her grapes are ripe. And the angel thrust in his sickle on the earth, and cut down the grapes of the vineyard of the Earth, and cast them into the great wine vat of the wrath of God: and the wine vat was trodden without the city. And the blood came out of the vat, even unto the horse bridles, by the space of a thousand and six hundred furlongs.

\index{Antichrist}\index{Hebrews, Epistle to the}Now he proceeds to describe God’s judgment, especially against the Antichristians and all the ungodly. This section could be joined with the following material, as it deals with the same argument. It pertains to the consolation and confirmation of the faithful, persecuted by Antichrist. Some people think there will never be any judgment. However, because they oppress their neighbors, they believe they will never feel God’s displeasure. Moreover, the faithful are tempted when they see the wicked prosper and themselves decline daily. Therefore, they think the Lord is taking too long. They expostulate with the Lord, saying, “When will the injustices end?” If Christ will come in judgment, why does he delay it, and cause such great suffering to his people? The Lord now shows that the judgment will certainly come, and will occur when all things are ripe—that is, when the iniquities of the Amorites are complete and the measure of iniquity is full. When wicked people are ripe, the Lord will come to judge. In the meantime, we must persevere in constancy and patience, like farmers waiting for the harvest and the grape harvest. If anyone rebels through impatience, he will not be approved by the Lord, as the Apostle states from the prophet in Hebrews 10. And just as we may desire and long for the harvest and the grape harvest, so we ought not to expostulate with God because he tarries longer than we wish; likewise, we ought not to contend with him as to why he comes so late in judgment. And just as the harvest and the grape harvest are certainly expected and arrive, so without a doubt God will punish the wicked and save the godly. These are truly tastes of what is to follow abundantly and are more expressly declared, and are appended to the former matters, because they pertain to the consolation of the godly.

And to the intent that all things might be more clear, using parables to present them before our eyes. He employs two parables drawn from the prophets and the doctrine of the Gospel. For the prophets frequently depict God’s judgment through harvest and vintage. Certainly in Joel 3, the Lord says: “I will sit in the valley of Josaphat to judge all nations. Put in the sickle, for the harvest is ripe,” and so forth. It is also well known what is said on the same matter in the story of the Gospel. We must therefore write these things more deeply into our hearts, and fear God, and persevere in his redemption with patience.

\index{Jerome, Saint}\index{Daniel (Book of)}First is treated the parable of harvest, then the parable of vintage: both showing that the Lord will be judge and that in his appointed time, against all those who either think there shall be no judgment, or question the Lord’s timing, saying he comes too slowly and late, \&c. And first is described the owner of the harvest, the Lord himself and judge Jesus Christ. He is said to be like unto the son of man: not because he is not now the very son of man, nor because he shall not come to judgment in the very human nature which he once took from us and never laid it aside (for he is truly the son of man, and remains on the right hand of the Father, and shall truly come in human nature to judge the living and the dead), but he seems to have alluded to Daniel, and to have expressed his speech, saying: “I looked in a mighty vision, and lo, there came one in the clouds as it were the son of man,” \&c. Where we read also the description of the judgment against the beast. And therefore he has here made mention also of a cloud: “and I saw a white cloud, and one sitting on the cloud,” \&c. Moreover, the angels in the Acts say, “so shall he come, as they have seen him go up into heaven.” And they saw him taken up, and a cloud to receive him and convey him out of their sight. Therefore shall he come again in a cloud unto judgment. We read oftentimes in the Psalms that God sits on a white cloud. By the way, therefore, is signified the deity of the judge. Therefore this judge is very God and very man, the Savior of the faithful, the revenger and judge of the unbelievers. We are sent therefore by Jerome to the seventh chapter of Daniel.

Then he wears a golden crown on his head: not that there is any corruptible gold in Heaven, but for corruptible men so he speaks, that they may understand their judge to be the high king; and may gather thereof, that no man is able to resist the power of this king. For otherwise our Lord hath no need of any corruptible gold.

Finally, our Lord here has a sickle, and that right sharp. Whereby is signified his judgment exceeding straight, and destruction of the wicked. In the third chapter of Matthew, the judgment of the Lord is compared to a fan, by blessed John. He adds that the axe is laid at the root of the tree: whereby he signified that certain judgment was at hand, or rather destruction.

Now follows an exposition of the proceeding of the judgment, and he perseveres in the parable. For he speaks as if a servant returning home out of the fields, did show unto his Master who looked for the hour of harvest, that the corn was now ripe (the hardness of the grain is a token of ripeness) and that it is time to be reaped. For else it is no need to admonish him that knows all things of any thing, that he remembers not: much less of the hour of judgment which no angel knows, but the Father alone. Therefore we ascribe this wholly to the parable: and we understand that a certain hour of judgment is appointed, which when it shall come, the godly will be delivered without delay, and the ungodly condemned. Another angel, says he, came forth. For before we heard how diverse came forth. This cries with a loud voice, as one that will tell of a matter most great and certain, and to be declared in the church with exceeding great outcry, to the comfort of the faithful, who ought nothing to doubt of the judgment, and to the terror of the wicked, who seem to contain the same. And this crying angel comes out of the temple. For we heard before that John saw a temple in heaven. And where the crier of the judgment comes out of the temple, it signifies that no unrighteousness of the judge is to be imagined. For the temple is consecrated to holiness and righteousness, and is called the house of God. Justly therefore he judges, and in just time he judges, and justly executes all things. The angel bids the judge do that which he of himself was about to do. Thrust in the sickle, says he, and reap. Two causes are alleged. First, for the hour is come that you should reap. Therefore a certain hour of judgment is appointed, which when it comes, the judgment shall be most certainly. And it is come for you, says he, for all judgment is given to the Son. Then, for the corn of the earth is ripe. As though he should say: the iniquity of earthly men is grown up to the highest, therefore it is reasonable that it should be cut down. And God alone knows when the iniquity of the earth is fulfilled. But when it shall come thereunto, there shall need no great preparation, deciding, or pondering of causes. At one word he finishes the judgment, and the execution of the same, and as it were swallows up and devours the whole earth in a moment, saying: herewith he thrust in his sickle, which sat upon the cloud, on the earth, and the earth was reaped. The rest of the things which seem to belong hereunto, take out of chapter 13 of Matthew. And that which he has said hitherto, he repeats, and beats in by another parable. For by this he shadows the same which the other parable did commend. That plenty makes for the plainer evidence, and beats in most diligently the certainty and verity of the judgment, lest herein we should doubt anything, and waver with the unfaithful world. The parable is taken from vintage. The same is used very often by the prophets, speaking of the destruction of any nation. And the Lord also in the gospel compares his people to a vine. And the angel holds in his hand a sharp sickle. He represents a figure of Christ, who has all power of judgment alone. A sharp sickle is the straight judgment, as was spoken of the sickle before. This angel comes out of the temple also, to wit a judge most righteous. Unto him cries another angel, which had power over fire, which comes out from the altar. For before we heard that there is an altar in the temple, and that under this altar do rest the souls of the blessed martyrs. Here therefore is figured that God does now remember the bloodshed of his servants, who for the profession of the only altar (that is Christ the priest and only sacrifice) were slain, and now to proceed to take vengeance hitherto long delayed. Therefore this angel is said to have power over fire. Fire many times in the Psalms signifies God’s vengeance. This angel therefore is here, as it were master of execution, and captain of vengeance. For angels in Daniel also, as God’s ministers, are said to have rule over things: not that we should worship and honor these ministers, but the Lord that works by them. The sun and moon are the lights of the world: but therefore no wise man will worship them. Here is signified plainly that vengeance is certainly prepared for them which shed innocent blood on the earth, and that this vengeance shall chiefly be executed in the end of this world. Albeit that he punishes nevertheless grievously before the end also here on earth, namely parricides: in so much that the Psalmograph says, men of blood shall not live half their time.

And as in the parable of harvest, harvest was finished with a short sentence: So here also, the vintage is ended with few words. For as soon as the ungodly shall see Christ in the clouds, with the prints of his wounds, and his Saints with him, whom they have scorned, hated, persecuted, and slain, they will immediately gather, that they by their just desert must be allotted with devils, whom they have followed and served. Therefore, there shall need no long discussion of the matter. Every man’s conscience shall accuse him, and the sins of every man shall be manifest to all creatures. The ungodly shall stand before the judge with great confusion, in utter contempt, in pain and fear, and sorrows not to be expressed, and shall go straight ways into pains and torments that shall never have end. Hereof I say, it behooves often to make mention, hereof it becomes often to warn all men, that they may beware in time, and take heed to themselves.

\index{John, Saint}However, Saint John himself describes the everlasting damnation and vengeance that God executes upon his enemies. And he uses the image of a wine press or wine vat, so that he may linger in the allegory, and it is made outside the city. By explanation, he calls it the great wine vat of God’s wrath. For it is hell, or the place of punishment and condemnation. Into this wine vat shall be gathered the clusters of the earth, or the grapes of the earth, that is, the earthly and ungodly men. And the city of God is heaven itself, the seat of the blessed, which shall afterward be described most abundantly in chapter 21. But that wine press is set outside the city. For in another place in the Gospel, the Lord also says that the wicked must be cast out into utter darkness, where there is weeping and gnashing of teeth.

But this wine vat is rightly called the wine vat of God’s wrath. For the wrath of God is executed therein: and those with whom God is angry for their sins are shut up therein, that there they may, according to their demerits, be tormented and vexed forever, and without end. And he calls it great, for that the place is wide enough to receive all the ungodly. As also Isaiah has admonished at the end of the thirtieth chapter. Others read of the great wrath of God.

It is added that from the fat or wine press no wine flows, but blood, and that in great abundance. This he illustrates with a marvelous and horrible hyperbole. The blood flowed far and wide, for the distance of a thousand six hundred furlongs. Again, it was very deep; for it rose to the bridles of the horses, those I mean, which went and wrestled in the blood, namely, in their own blood. By this hyperbolic speech is signified that the multitude of the ungodly will be greatest, and that God will most abundantly avenge that immeasurable blood which the wicked have spilled on earth. They delighted themselves while living on earth with wars, slaughter, persecutions, and martyrdoms; therefore, God, most just, will give them in another world enough blood, so that being drowned in their own blood up to their chins, they may seem to bathe themselves in their own blood. And here we must remember that the horses prepared for battle of which we spoke in chapter 9 will be drowned in everlasting torments. Thus, thus at the last will the Lord avenge himself upon his enemies. Let us call upon him, and abide patiently and valiantly. The Lord grant us his grace.
